\documentclass[11pt]{article}
\usepackage[a4paper,margin=28mm]{geometry}
\usepackage[T1]{fontenc}
\usepackage{lmodern,amsmath,amssymb,amsthm,mathtools,microtype}
\usepackage[hidelinks]{hyperref}
\newcommand{\R}{\mathbb R}
\newcommand{\Z}{\mathbb Z}
\newcommand{\Pp}{\mathbb P}
\newcommand{\E}{\mathbb E}
\renewcommand{\Pr}{\mathbb P\!\operatorname{r}}
\newtheorem{theorem}{Theorem}[section]
\newtheorem{lemma}[theorem]{Lemma}
\newtheorem{proposition}[theorem]{Proposition}
\numberwithin{equation}{section}
\title{Local well-posedness for transport-noise Navier--Stokes\protect\\
with $B^{-1/2}_{6,4}(\R^3)$ initial data}
\author{Kaiwen Jin\\\small\texttt{yui.ui.siki@gmail.com}}
\date{October 15, 2026}
\begin{document}
\maketitle
\begin{abstract}
We prove local well-posedness of the three-dimensional incompressible
Navier--Stokes equation with prescribed smooth, compactly supported,
divergence-free Ito transport noise and deterministic initial velocity in
the inhomogeneous space $B^{-1/2}_{6,4}(\R^3)$. Under strict stochastic
parabolicity, with no further smallness assumption, the solution is heat-mild
and belongs almost surely to $C_tB^{-1/2}_{6,4}\cap L^4_tL^6_x$ up to a
positive stopping time. Pathwise uniqueness holds in this entire class, and
initial-data dependence is continuous in probability in both norms on common
local intervals. The proof combines the published linear stochastic Stokes
maximal-regularity theorem of Agresti and Veraar with a cutoff depending on
the accumulated $L^4_tL^6_x$ norm. This gives a nonlinear local theory at the
negative-half endpoint without treating the derivative noise as a small
perturbation.
\end{abstract}

\section{Introduction and main result}
The transport-noise Navier--Stokes equation replaces part of the advecting
velocity by random spatial fields. After projecting onto incompressible
velocities, the Ito equation on $\R^3$ takes the form
\begin{equation}\label{eq:NS}
 du=\big[\nu\Delta u-\Pp\operatorname{div}(u\otimes u)\big]dt
       +\sum_{k=1}^K\Pp[(b_k\cdot\nabla)u]\,dW_t^k.
\end{equation}
Here $\nu>0$ is viscosity, $\Pp$ is the Helmholtz projection, and $W^k$
are independent standard real Brownian motions. We consider real fields
$b_k\in C_c^\infty(\R^3;\R^3)$ with $\operatorname{div}b_k=0$ and assume
that, for some $\delta>0$,
\begin{equation}\label{eq:parabolic}
 \nu I-\frac12\sum_{k=1}^Kb_k(x)\otimes b_k(x)\geq\delta I
 \qquad (x\in\R^3).
\end{equation}
This is the strict stochastic parabolicity condition. It permits the
first-order noise to enter the principal linear equation, with constants
depending on the prescribed fields.

Rough initial velocities require a solution theory compatible with viscous
smoothing. The Navier--Stokes scaling assigns spatial smoothness $-1+3/q$
at integrability $q$. At $q=6$ this is $-1/2$, while the corresponding
space-time norm is $L^4_tL^6_x$. These are high-frequency critical indices;
the Besov space used here is inhomogeneous.
Agresti and Veraar developed stochastic Stokes maximal regularity and
critical-space nonlinear theory for transport-noise equations
\cite{AV24}. Their nonlinear theorem covers spatial integrability below
six at the lowest smoothness. The endpoint is part of the rough-data
question discussed in \cite[Theorem 8.26 and Open Problem 10]{AV25}.
The linear theorem itself allows the Sobolev shift needed at $q=6$.
Its nonlinear conclusion cannot simply be evaluated there: the critical
time weight reaches the excluded boundary of the stochastic
maximal-regularity range. We use that linear theorem and a separate
heat-mild quadratic estimate to obtain the endpoint.

Write $X=B^{-1/2}_{6,4,\sigma}(\R^3)$, with $\sigma$ denoting divergence-free
distributions. Work on any complete right-continuous filtered probability
space carrying the Brownian motions above. By an adapted heat-mild solution
on $[0,\tau]$ we mean an adapted continuous $X$ process, with a progressive
$L^6$ representative at almost every positive time, satisfying
\begin{align}\label{eq:mild}
 u(t)={}&e^{\nu t\Delta}u_0
 -\int_0^te^{\nu(t-s)\Delta}\Pp\operatorname{div}(u\otimes u)(s)\,ds\notag\\
 &+\sum_{k=1}^K\int_0^te^{\nu(t-s)\Delta}
                  \Pp[(b_k\cdot\nabla)u(s)]\,dW_s^k,
 \qquad 0\leq t\leq\tau.
\end{align}
The identity is in distributions, equivalently in $H^{-2,6}$, with the
continuous versions of the convolutions. Stochastic integrals are first
defined at deterministic times after localization and then restricted to the
stopped interval. This specifies the heat semigroup in both convolutions.

\begin{theorem}\label{thm:main}
Let $\nu>0$ and let the finite family $(b_k)_{k=1}^K$ satisfy the assumptions
above, including \eqref{eq:parabolic}. For every deterministic $u_0\in X$
there exist a stopping time $0<\tau\leq1$ almost surely and an adapted
heat-mild solution of \eqref{eq:NS} such that
\[
 u\in C([0,\tau];X)\cap L^4(0,\tau;L^6(\R^3))
 \quad\text{almost surely}.
\]
Any two heat-mild solutions in this class with the same initial value and
Brownian motions agree on their common lifetime.

If deterministic $u_{0,n}\to u_0$ in $X$, the solutions can be chosen on the
same Brownian space with positive lifetimes $\tau_n,\tau\leq1$ so that
\[
 \|u_n-u\|_{C([0,\tau_n\wedge\tau];X)}
 +\|u_n-u\|_{L^4(0,\tau_n\wedge\tau;L^6)}\longrightarrow0
 \quad\text{in probability}.
\]
There is also a stopping time $0<\rho\leq\tau$ almost surely with
$\Pr(\tau_n<\rho)\to0$. No extra smallness of $u_0$ or of the fields is
assumed.
\end{theorem}

The result addresses the specified $B^{-1/2}_{6,4}$ case of the rough-data
question, for all the prescribed compactly supported fields satisfying
\eqref{eq:parabolic}. It does not assert global existence or determine the
largest possible critical initial space.

The proof treats the full linear derivative noise before the quadratic
term. The published linear estimate yields a bounded causal inverse for
the identity minus the heat stochastic-convolution operator on
$L^4(\Omega;L^4_tL^6_x)$. A cutoff depending on the past accumulated
$L^4_tL^6_x$ norm then makes the drift globally Lipschitz with an arbitrarily
small constant. Stopping before the cutoff acts recovers the original
equation. For uniqueness among solutions with only almost-sure norms,
we use bounded-norm random slabs and the stopped linear estimate.

Section~\ref{sec:estimates} proves the heat estimates and specifies the
published linear input, including the stopping-time inverse.
Section~\ref{sec:proof} constructs the local solution and proves uniqueness
and stability. We use $\widehat f(\xi)=\int e^{-ix\cdot\xi}f(x)\,dx$;
$\Pp$ has symbol $I-\xi\otimes\xi/|\xi|^2$. All Besov and Sobolev spaces
are inhomogeneous, and $\|\cdot\|_q$ denotes the spatial $L^q$ norm.
\section{Heat estimates and the linear stochastic equation}\label{sec:estimates}

All spaces below are on $\R^3$, with vector or tensor values as appropriate.
Write $J=(1-\Delta)^{1/2}$ and
$H^{s,q}=\{f:J^sf\in L^q\}$. Fix an inhomogeneous smooth dyadic
partition $(\Delta_j)_{j\geq-1}$, with one low-frequency block and annuli
of size $2^j$ for $j\geq0$. We use the norm
\begin{equation}\label{eq:Xnorm}
 \|f\|_X=\Big(\sum_{j\geq-1}2^{-2j}\|\Delta_jf\|_6^4\Big)^{1/4},
 \qquad X=B^{-1/2}_{6,4,\sigma}.
\end{equation}
The subscript $\sigma$ denotes the distributional kernel of divergence.
Different admissible partitions give equivalent norms by finite overlap and
uniform $L^1$ bounds for their convolution kernels. The space $X$ is a
separable Banach space and embeds continuously into $H^{-2,6}_\sigma$.
Indeed, if $a_j=2^{-j/2}\|\Delta_jf\|_6$, the high-frequency contribution
to the latter norm is bounded by
$C\sum_{j\geq0}2^{-3j/2}a_j\leq C\|(a_j)\|_{\ell^4}$.
Completeness follows by compatible dyadic block limits, and separability by
the closed embedding into $\ell^4(L^6)$.

For $0<T<\infty$ set
\[
 E_T=L^4(0,T;L^6_\sigma),\qquad
 \mathcal E_T=L^4(\Omega;E_T)_{\mathrm{prog}}.
\]
Here the second space consists of progressively measurable processes. It is
a closed Banach subspace of the indicated Bochner space. For a stopping time
$\gamma\leq T$, $\|v\|_{E_\gamma}$ means
$(\int_0^\gamma\|v(s)\|_6^4\,ds)^{1/4}$. Constants can depend on $T$,
$\nu$, the fixed fields and the parabolicity margin; no uniformity over
coefficient families is asserted.

\subsection{Deterministic heat estimates}
Let $S(t)=e^{\nu t\Delta}$ and, for tensor-valued $F$, define
\[
 (DF)(t)=\int_0^tS(t-s)\Pp\operatorname{div}F(s)\,ds.
\]
The following estimates include the time-continuous trace that will be
needed for the stopping construction.

\begin{lemma}\label{lem:heat}
For each finite $T$ there are finite constants such that
\begin{align}
 \|DF\|_{E_T}&\leq C_D\|F\|_{L^2(0,T;L^3)},\label{eq:D-E}\\
 \|DF\|_{C([0,T];X)}&\leq C_X(T)\|F\|_{L^2(0,T;L^3)},
 \qquad DF(0)=0,\label{eq:D-X}\\
 \|S(\cdot)f\|_{E_T}+\|S(\cdot)f\|_{C([0,T];X)}
 &\leq C_0(T)\|f\|_X.\label{eq:initial}
\end{align}
Moreover $\Pp\operatorname{div}:L^3\to H^{-2,6}_\sigma$ is bounded.
The constant $C_D$ can be chosen independently of $T$.
\end{lemma}

\begin{proof}
The Gaussian derivative bound, Young's inequality and boundedness of $\Pp$
on $L^6$ give
\begin{equation}\label{eq:heatkernel}
 \|S(t)\Pp\operatorname{div}F\|_6\leq Ct^{-3/4}\|F\|_3.
\end{equation}
The derivative contributes $t^{-1/2}$ and the $L^3$ to $L^6$ estimate
contributes $t^{-1/4}$. Extend $\|F(t)\|_3$ by zero outside $(0,T)$.
The one-dimensional Hardy--Littlewood--Sobolev inequality with fractional
order $1/4$ maps $L^2(\R)$ to $L^4(\R)$ and proves \eqref{eq:D-E}.
We use the finite-exponent forms of these inequalities; see
\cite[Corollary 6.3 and Theorem 6.4]{Tao2}.

For the trace, choose enlarged blocks $\widetilde\Delta_j$ equal to one on
the support of $\Delta_j$. The localized heat kernel satisfies, for $j\geq0$,
\[
 \|\Delta_jDF(t)\|_6\leq C2^{3j/2}
 \int_0^t e^{-c\nu2^{2j}(t-s)}
          \|\widetilde\Delta_jF(s)\|_3\,ds.
\]
To obtain this bound, scale the multiplier to a fixed annulus. Its smooth
kernel has $L^{6/5}$ norm at most
$C2^{3j/2}e^{-c\nu2^{2j}(t-s)}$: the derivative costs $2^j$, the spatial
integrability change costs $2^{j/2}$, and integration by parts absorbs the
polynomial factors in $\nu2^{2j}(t-s)$ into the exponential.
Cauchy--Schwarz in time therefore yields
\begin{equation}\label{eq:blocktrace}
 \sup_{t\leq T}2^{-j/2}\|\Delta_jDF(t)\|_6
 \leq C\|\widetilde\Delta_jF\|_{L^2(0,T;L^3)}.
\end{equation}
The low block is bounded by
$C\sqrt T\|\widetilde\Delta_{-1}F\|_{L^2L^3}$, using its fixed Fourier
support and heat contraction. For the high blocks, mixed-norm interchange
and the upper Littlewood--Paley estimate give
\begin{align*}
 \Big\|\big(\|\widetilde\Delta_jF\|_{L^2L^3}\big)_{j\geq0}\Big\|_{\ell^4}
 &\leq\Big\|\big(\|\widetilde\Delta_jF\|_{L^2L^3}\big)_{j\geq0}\Big\|_{\ell^3}\\
 &\leq\left\|\Big(\sum_{j\geq0}|\widetilde\Delta_jF|^3\Big)^{1/3}
        \right\|_{L^2L^3}\\
 &\leq\left\|\Big(\sum_{j\geq0}|\widetilde\Delta_jF|^2\Big)^{1/2}
        \right\|_{L^2L^3}
 \leq C\|F\|_{L^2L^3}.
\end{align*}
The mixed-norm direction uses $3\geq2$; the last estimate is the upper
square-function bound at spatial exponent $3$
\cite[Proposition 4.1]{Tao2}, \cite[Proposition 5.3]{Tao4}.
Together with \eqref{eq:blocktrace}, this proves the uniform $X$ bound.
Approximation by smooth functions in $L^2L^3$ gives continuity into $X$ and
zero initial trace, proving \eqref{eq:D-X}.

For \eqref{eq:initial}, first let $T=1$. The low block is in $L^6$.
For the high blocks put $a_j=2^{-j/2}\|\Delta_jf\|_6$.
On $2^{-2(n+1)}<t\leq2^{-2n}$, $n\geq0$, the annular heat bound implies
\[
 2^{-n/2}\left\|S(t)\sum_{j\geq0}\Delta_jf\right\|_6
 \leq C\sum_{j\geq0}2^{(j-n)/2}
                  e^{-c\nu2^{2(j-n-1)}}a_j.
\]
The sequence $2^{\ell/2}e^{-c\nu2^{2(\ell-1)}}$ lies in $\ell^1(\Z)$.
Integration over these time intervals and discrete Young's inequality give
the $E_1$ bound. For continuity in $X$ at zero, each
$(S(t)-I)\Delta_jf$ tends to zero in $L^6$ and is bounded by
$2\|\Delta_jf\|_6$. Dominated convergence in \eqref{eq:Xnorm} proves the
claim. Commutation with the blocks gives the uniform $X$ bound and continuity
at later times. If $T>1$, heat contraction after time $1$ gives finite
$T$-dependent constants.

Finally the heat representation of $J^{-2}$ and \eqref{eq:heatkernel} with
viscosity one give
\[
 \|\Pp\operatorname{div}F\|_{H^{-2,6}}
 \leq C\int_0^\infty e^{-s}s^{-3/4}\,ds\,\|F\|_3
 \leq C\|F\|_3.
\]
\end{proof}

\subsection{The published linear estimate}
The tensor convention is $(v\otimes b)_{ij}=v_i b_j$ and
$(\operatorname{div}F)_i=\sum_j\partial_jF_{ij}$. Thus $\operatorname{div}b_k=0$
implies
\begin{equation}\label{eq:Bbound}
 B_kv:=\Pp[(b_k\cdot\nabla)v]=\Pp\operatorname{div}(v\otimes b_k),
 \qquad
 \|Bv\|_{H^{-1,6}(\ell^2)}\leq C_b\|v\|_6.
\end{equation}
Here $Bv=(B_kv)_{k=1}^K$, and the bound follows from the bounded multiplier
$J^{-1}\partial_j$ on $L^6$ and boundedness of the fields.
The projection remains outside the variable coefficient derivative.

\begin{proposition}[Agresti--Veraar, specialized]\label{prop:linear}
For every finite $T$ and progressive
$g\in L^4(\Omega;L^4(0,T;H^{-1,6}_\sigma(\ell^2)))$, the equation
\begin{equation}\label{eq:linear}
 dz=\nu\Delta z\,dt+(Bz+g)\,dW,\qquad z(0)=0,
\end{equation}
has a unique strong solution with
\begin{equation}\label{eq:linearbound}
 \|z\|_{\mathcal E_T}
 +\|z\|_{L^4(\Omega;C([0,T];X))}
 \leq C_L(T)\|g\|_{L^4(\Omega;L^4(0,T;H^{-1,6}(\ell^2)))}.
\end{equation}
The same assertion holds with $B=0$. The strong identity is in $H^{-2,6}$.
It admits stopping-time restriction and pathwise-integrability localization;
in particular, two local strong solutions with the same forcing and zero
initial value agree before a common stopping time whenever their $E$ norms
are almost surely finite.
\end{proposition}

\begin{proof}[Source and parameter verification]
Use \cite[Assumption 3.1, Theorem 3.2 and Remark 3.3]{AV24} with
\[
 (p,q,\kappa,\eta)=(4,6,0,-1),\qquad
 X_0=H^{-2,6}_\sigma,\quad X_1=L^6_\sigma,
 \quad X_{1/2}=H^{-1,6}_\sigma.
\]
Here $\eta$ denotes the spatial shift called $\delta$ in that source;
it is unrelated to the positive margin $\delta$ in \eqref{eq:parabolic}.
The source linear assumption allows every real spatial shift. Its trace is
\[
 (X_0,X_1)_{3/4,4}=B^{-1/2}_{6,4,\sigma}=X,
 \qquad 1+\eta-2(1+\kappa)/p=-1/2.
\]
Take $a^{ij}=\nu\delta_{ij}$ and the source coefficients $c=h=0$.
The additional pressure-dependent drift in its general Stokes system then
vanishes. The noise operator is exactly $\Pp(b_k\cdot\nabla)$.
The source coefficient Sobolev regularity may be chosen of order $2$ and
integrability $6$: $2>1=-\eta$ and $2>3/6$. Smooth compactly supported
$b_k$ satisfy this condition. For the whole-space version, Remark 3.3
permits a constant background at infinity: $a=\nu I$ is that background
and $b=0$ outside a compact set. No global $L^6$ integrability of the
nonzero constant matrix is needed. Its strict stochastic parabolicity
condition is precisely \eqref{eq:parabolic}. Extend the finite noise family
by zeros in $\ell^2$.

Theorem 3.2 now gives \eqref{eq:linearbound}, including its continuous trace.
The stopped and localized statement is the corresponding stochastic maximal
regularity assertion: the proof of that theorem establishes the
$\mathrm{SMR}^{\bullet}_{p,\kappa}$ property, and its Lemma 3.7 gives the
stopped a priori estimate. Apply the pathwise-data localization in
\cite[Definitions 3.7--3.8 and Proposition 3.11]{AV25} to this couple.
This gives the stated local strong uniqueness as well as the deterministic
interval estimate. These are linear results; no nonlinear endpoint theorem
is being imported.
\end{proof}

\subsection{Heat stochastic convolutions and a causal inverse}
For $v\in\mathcal E_T$ define, initially at deterministic times,
\begin{equation}\label{eq:K}
 (Kv)(t)=\sum_{k=1}^K\int_0^tS(t-s)B_kv(s)\,dW_s^k.
\end{equation}
The integral is well defined in $H^{-1,6}$, since \eqref{eq:Bbound} and the
finite-horizon embedding $L^4_t\subset L^2_t$ apply. The usual Banach-space
Ito integration on these Sobolev spaces is available; see
\cite[Section 2.5 and Proposition 2.4]{AV25}.
The heat-only case of Proposition~\ref{prop:linear} shows that $K$ has a
continuous $X$ version and
\begin{equation}\label{eq:Kbound}
 \|Kv\|_{\mathcal E_T}+\|Kv\|_{L^4(\Omega;C([0,T];X))}
 \leq C_K(T)\|v\|_{\mathcal E_T},\qquad Kv(0)=0.
\end{equation}

For clarity, the strong and heat-mild identities agree at this regularity.
If $Z$ solves the heat-only equation with forcing $g=Bv$, pair it, at time
$s$, with $S(t-s)\phi$ for a Schwartz field $\phi$ and deterministic $t$.
The derivative of the test field is $-\nu\Delta S(t-s)\phi$, so the
Laplacian pairings cancel and
\begin{equation}\label{eq:test}
 \langle Z(t),\phi\rangle
 =\sum_k\int_0^t\langle B_kv(s),S(t-s)\phi\rangle\,dW_s^k.
\end{equation}
This testing rule follows from deterministic time partitions and summation
by parts. Drift errors converge in $L^1_tH^{-2,6}$ and stochastic errors by
the scalar Ito isometry after localization. Thus \eqref{eq:test} identifies
$Z$ with \eqref{eq:K}, proving \eqref{eq:Kbound}. A countable dense set of
tests and rational times identifies their continuous versions. The same
argument for \eqref{eq:linear} gives $z=Kz+K_g$, with $K_g$ the heat
stochastic convolution of $g$.

These identities also hold after stopping the cumulative $E$ norm, for
processes whose $E$ norm is only almost surely finite.
The deterministic terms satisfy the corresponding strong identities in
$H^{-2,6}$: $\Delta S(\cdot)f\in L^4_tH^{-2,6}$ and, by
Lemma~\ref{lem:heat}, $\Delta DF\in L^4_tH^{-2,6}$ and
$\Pp\operatorname{div}F\in L^2_tH^{-2,6}$.
Consequently a heat-mild solution in $E\cap C X$ is a strong
$H^{-2,6}$ solution. The nonlinear drift is only required to be
$L^1_tH^{-2,6}$; it is not used as $L^4$ drift forcing in
Proposition~\ref{prop:linear}. Evaluation of the continuous convolution at
a stopping time means restriction of this version, not an Ito integral
with a random terminal-time kernel.

\begin{lemma}\label{lem:inverse}
On $\mathcal E_T$, $I-K$ has a bounded linear causal inverse $R_T$.
There is a finite $M_T\geq1$ such that, for every stopping time
$\gamma\leq T$ and progressive $w,v$ satisfying $w-Kw=v$ on $[0,\gamma]$
with finite $L^4(\Omega;E_\gamma)$ norms,
\begin{equation}\label{eq:stoppedinverse}
 \|w\|_{L^4(\Omega;E_\gamma)}
 \leq M_T\|v\|_{L^4(\Omega;E_\gamma)}.
\end{equation}
\end{lemma}

\begin{proof}
Given $v\in\mathcal E_T$, solve \eqref{eq:linear} with $g=Bv$ and put
$R_Tv=v+z$. Then
$\|R_Tv\|_{\mathcal E_T}\leq(1+C_L(T)C_b)\|v\|_{\mathcal E_T}$.
The identity $z=Kz+Kv$ gives $(I-K)R_Tv=v$.
Conversely, if $w-Kw=v$, then $z=w-v=Kw$ is the zero-initial strong solution
of \eqref{eq:linear} with forcing $Bv$, because $Bz+Bv=Bw$.
Linear uniqueness gives $w=R_Tv$. Equal inputs through a stopping time give
equal forcing there; local linear uniqueness proves causality.

For the stopped estimate, $z=w-v=Kw$ solves \eqref{eq:linear} up to
$\gamma$ with $g=Bv$. Extend \emph{this forcing} by zero after $\gamma$
and solve the linear equation on $[0,T]$. Local uniqueness in
Proposition~\ref{prop:linear} identifies its restriction with $z$.
Restrict \eqref{eq:linearbound} and add $v$ to obtain
\eqref{eq:stoppedinverse} with $M_T=1+C_L(T)C_b$.
The solution itself is not extended by zero, so no jump at $\gamma$ is
introduced. This construction does not require $\|K\|<1$ or a bounded
pathwise inverse.
\end{proof}

\section{Local solutions, uniqueness and stability}\label{sec:proof}

\subsection{A cutoff depending on the past}
For a path $u\in E_T$ let
\[
 A_u(t)=\Big(\int_0^t\|u(s)\|_6^4\,ds\Big)^{1/4}.
\]
It is continuous, nondecreasing, starts at zero, and satisfies
$|A_u(t)-A_v(t)|\leq\|u-v\|_{E_t}$. For a progressive process it is
adapted. Fix a nonincreasing Lipschitz function $\theta:[0,\infty)\to[0,1]$
with $\theta=1$ on $[0,1]$ and $\theta=0$ on $[2,\infty)$. Write
$L_\theta$ for its Lipschitz constant. To truncate the quadratic drift
without using future values, set
\[
 F_r(u)(t)=\theta(A_u(t)/r)\,u(t)\otimes u(t),\qquad r>0.
\]

\begin{lemma}\label{lem:cutoff}
For all $u,v\in E_T$,
\begin{align}
 \|F_r(u)\|_{L^2L^3}&\leq4r^2,\label{eq:cutoffbound}\\
 \|F_r(u)-F_r(v)\|_{L^2L^3}
 &\leq C_\theta r\|u-v\|_{E_T},
 \qquad C_\theta=\max\{8,6+9L_\theta\}.\label{eq:cutofflip}
\end{align}
These bounds are pathwise, and $F_r$ preserves progressive measurability.
\end{lemma}

\begin{proof}
The support of the cutoff lies before $A_u$ reaches $2r$, where the
$E$ norm is at most $2r$; Holder's inequality gives
\eqref{eq:cutoffbound}. Put $d=\|u-v\|_{E_T}$.
If $d\geq r$, the sum of the two bounds is $8r^2\leq8rd$.
If $d<r$, then on the union $J_*$ of the cutoff supports both cumulative
norms are less than $3r$. Hence
$\|\mathbf1_{J_*}u\|_{E_T},\|\mathbf1_{J_*}v\|_{E_T}\leq3r$.
With $\theta_u(t)=\theta(A_u(t)/r)$, decompose
\[
 F_r(u)-F_r(v)
 =\theta_u(u\otimes u-v\otimes v)
   +(\theta_u-\theta_v)v\otimes v.
\]
The two $L^2L^3$ bounds are $6rd$ and $9L_\theta rd$, respectively.
Outside $J_*$ the difference is zero. This proves \eqref{eq:cutofflip}.
The measurability claim follows from adaptedness and continuity of $A_u$.
\end{proof}

\subsection{Construction and positive lifetime}
Use $T=1$ and abbreviate $M=M_1$ and $C_0=C_0(1)$. Enlarge the positive
constants if necessary and choose
\begin{equation}\label{eq:radius}
 r=(4MC_DC_\theta)^{-1}.
\end{equation}
For deterministic $f\in X$ consider the map on the whole space
$\mathcal E_1$
\[
 \Phi_f(u)=R_1\big[S(\cdot)f-DF_r(u)\big].
\]
Lemmas~\ref{lem:heat}, \ref{lem:inverse} and \ref{lem:cutoff} imply
\[
 \|\Phi_f(u)-\Phi_f(v)\|_{\mathcal E_1}
 \leq MC_DC_\theta r\|u-v\|_{\mathcal E_1}
 =\tfrac14\|u-v\|_{\mathcal E_1}.
\]
The map takes $\mathcal E_1$ into itself. In particular, the pathwise bound
\eqref{eq:cutoffbound} gives the necessary fourth probability moment of
the truncated forcing. Picard iteration converges in this Banach space to
a unique fixed point $U_f$, for every size of $f$. Its equation is
\begin{equation}\label{eq:truncated}
 U_f=S(\cdot)f-DF_r(U_f)+KU_f.
\end{equation}
Equations~\eqref{eq:D-X}, \eqref{eq:initial} and \eqref{eq:Kbound} supply
an adapted continuous $X$ version with $U_f(0)=f$. It agrees with the
progressive $L^6$ representative almost everywhere. Separability of $X$
and the usual completion of the filtration permit one common null-set
modification for the continuous identity.

Define
\begin{equation}\label{eq:lifetime}
 \tau_f=\inf\{t\in[0,1]:A_{U_f}(t)\geq r\}\wedge1,
 \qquad \inf\varnothing=1.
\end{equation}
Continuity and adaptedness make this a stopping time, and
$A_{U_f}(0)=0$ implies $\tau_f>0$ almost surely.
On $[0,\tau_f]$ the cutoff equals one. Thus
$u_f=U_f|_{[0,\tau_f]}$ solves \eqref{eq:mild} with paths in
$C([0,\tau_f];X)\cap E_{\tau_f}$. Only positive-time $L^6$ values are
used in the quadratic term; no $L^6$ initial datum is required.

\subsection{Uniqueness in the whole stated class}
Let $u,v$ be any two solutions in the class of Theorem~\ref{thm:main},
with the same initial value and Brownian motions. Fix an integer $m\geq1$
and put $\gamma=\tau_u\wedge\tau_v\wedge m$.
Set $\varepsilon=(4M_mC_D)^{-1}$. Starting at $s_0=0$, define
\[
 s_{j+1}=\inf\left\{t\geq s_j:
   \int_{s_j}^t(\|u(a)\|_6+\|v(a)\|_6)^4\,da
                      \geq\varepsilon^4\right\}\wedge\gamma.
\]
The norm integrands are zero after $\gamma$ for this definition.
These are stopping times because the accumulated adapted integrals are
continuous. Every completed threshold consumes $\varepsilon^4$ of a finite
pathwise integral, so almost every path reaches $\gamma$ in finitely many
slabs.

Assume inductively that $u=v$ through $s_j$. On $[0,s_{j+1}]$ put $w=u-v$;
it is zero before $s_j$. Extend the tensor integrand
\[
 F=\mathbf1_{(s_j,s_{j+1}]}(u\otimes u-v\otimes v)
\]
by zero outside its slab and put $v_0=-DF$. Subtracting the mild equations
gives $w-Kw=v_0$ up to $s_{j+1}$. By \eqref{eq:D-E} and Holder,
\begin{equation}\label{eq:slabbound}
 \|v_0\|_{E_{s_{j+1}}}
 \leq2C_D\varepsilon\|w\|_{E_{s_{j+1}}}.
\end{equation}
The stopped norm of $w$ is at most $2\varepsilon$, since its past is zero.
Thus $w$ and $v_0$ have finite fourth probability moments on this stopped
interval, although the original solutions have no assumed moments.
Apply \eqref{eq:stoppedinverse} and then \eqref{eq:slabbound}:
\[
 \|w\|_{L^4(\Omega;E_{s_{j+1}})}
 \leq2M_mC_D\varepsilon\|w\|_{L^4(\Omega;E_{s_{j+1}})}
 =\tfrac12\|w\|_{L^4(\Omega;E_{s_{j+1}})}.
\]
Hence $w=0$ almost everywhere, and continuity in $X$ gives equality at
every time through $s_{j+1}$. A countable intersection over $j$ and then
$m$ proves equality on the entire common lifetime. The estimate starts at
time zero with an already vanishing past; it uses no random restart theorem.

\subsection{Continuous dependence on common local intervals}
The same $\theta,r$ work for every deterministic datum. The fixed-point
estimate gives
\begin{equation}\label{eq:Estability}
 \|U_f-U_g\|_{\mathcal E_1}
 \leq MC_0\|f-g\|_X+\tfrac14\|U_f-U_g\|_{\mathcal E_1}.
\end{equation}
In particular, absorbing the last term yields
$\|U_f-U_g\|_{\mathcal E_1}\leq2MC_0\|f-g\|_X$.
Subtracting \eqref{eq:truncated} and using the continuous-trace estimates
for $S,D,K$ and \eqref{eq:cutofflip} then gives
\begin{equation}\label{eq:Xstability}
 \|U_f-U_g\|_{L^4(\Omega;C([0,1];X))}\leq C\|f-g\|_X.
\end{equation}

Suppose $f_n\to f$ in $X$. Restriction of these estimates to
$[0,\tau_{f_n}\wedge\tau_f]$ proves fourth-moment convergence in the sum
of the two norms asserted in the theorem, hence convergence in probability.
To show that the common intervals do not collapse, let
\[
 \rho=\inf\{t\in[0,1]:A_{U_f}(t)\geq r/2\}\wedge1.
\]
Then $0<\rho\leq\tau_f$ almost surely. If $\tau_{f_n}<\rho$, continuity
gives $A_{U_{f_n}}(\tau_{f_n})=r$ and
$A_{U_f}(\tau_{f_n})<r/2$. The reverse triangle inequality implies
$\|U_{f_n}-U_f\|_{E_1}\geq r/2$, and therefore
\[
 \Pr(\tau_{f_n}<\rho)
 \leq16r^{-4}\E\|U_{f_n}-U_f\|_{E_1}^4\longrightarrow0.
\]
This proves the stability assertion and completes the proof of
Theorem~\ref{thm:main}.

\section*{AI-use declaration}
OpenAI Codex was used to develop the argument, check source applications,
draft the text, and perform internal AI reviews and revisions. These reviews
do not constitute independent human verification.

\begin{thebibliography}{9}
\bibitem{AV24}
A. Agresti and M. Veraar,
\emph{Stochastic Navier--Stokes equations for turbulent flows in critical spaces},
Communications in Mathematical Physics \textbf{405}, article 43 (2024).
\href{https://doi.org/10.1007/s00220-023-04867-7}{doi:10.1007/s00220-023-04867-7}.

\bibitem{AV25}
A. Agresti and M. Veraar,
\emph{Nonlinear SPDEs and maximal regularity: an extended survey},
Nonlinear Differential Equations and Applications NoDEA \textbf{32},
article 123 (2025).
\href{https://doi.org/10.1007/s00030-025-01090-2}{doi:10.1007/s00030-025-01090-2}.

\bibitem{Tao2}
T. Tao, \emph{Lecture notes 2 for 247A}, UCLA (2006).
\url{https://www.math.ucla.edu/~tao/247a.1.06f/notes2.pdf}.

\bibitem{Tao4}
T. Tao, \emph{Lecture notes 4 for 247A}, UCLA (2006).
\url{https://www.math.ucla.edu/~tao/247a.1.06f/notes4.pdf}.
\end{thebibliography}
\end{document}
